\begingroup

\section{The original gamma convolution and the exact arithmetic coefficient maps}



\subsection{Source, attribution, and original coordinates}

This chapter proves the coefficient and cochain maps needed to integrate
\emph{Tau Gamma Convolution Descent}, whose original
\texttt{NOTE.tex} has SHA256
\begin{center}\footnotesize
\texttt{7cf984d87de405b815250b44454c4ee27c8ab46574da245b377ebc6e49327fa2}.
\end{center}
The analytic generating identities (GC.24)--(GC.26) were supplied by
the \emph{Zeta Function Literature \& Local Corpus Foundation} task on
13 September 2026. The proofs here verify and integrate that supplied
join; they do not claim its discovery. The original gamma generating
integral and monic family are the existing TG construction, with the
classical Meixner--Pollaczek parameters retained.

All scalar source coordinates remain
\[
 s=\tfrac12+it,\qquad S=s_1+\cdots+s_k=\tfrac k2+iu,
 \qquad u=t_1+\cdots+t_k.                                      \tag{GC.1}
\]
Here $k\ge1$ is an integer. We use inner products conjugate-linear in
the first argument. For a polynomial, an overline over its name means
coefficientwise conjugation. In particular
$\overline{P(c+iu)}=\overline P(c-iu)$ when $c,u$ are real.

\subsubsection{The analytic source and its actual theta boundary}

Let $D=-x\partial_x$. Let $V$ be the complex vector space of even
Schwartz functions $\phi$ on $\mathbb R$ with
$\phi(0)=0$ and $\int_{\mathbb R}\phi(x)\,dx=0$. Let $\mathscr B$ be
the smooth functions on $\mathbb R_{>0}$ such that for every
$a>0$ and every integer $j\ge0$,
\[
 D^jF(x)=O(x^a)\quad(x\downarrow0),\qquad
 D^jF(x)=O(x^{-a})\quad(x\to\infty).
\]
The original theta map and cochain grading are
\[
 \Theta:V\longrightarrow\mathscr B,\quad
 \Theta\phi(x)=\sum_{n\ne0}\phi(nx),\qquad
 [V\xrightarrow{\Theta}\mathscr B]\text{ in degrees }0,1.
                                                               \tag{GC.2}
\]
The sum and its differentiated sums converge at infinity by the
Schwartz bounds. Poisson summation gives its corresponding bounds at
zero: the zero Fourier coefficient vanishes by $\int\phi=0$, and
the original zero term vanishes by $\phi(0)=0$. Differentiation
preserves both conditions, because $D\phi(0)=0$ and integration by
parts gives $\int D\phi=\int\phi=0$. Thus $D\Theta=\Theta D$ on
these spaces. The rapid bounds make the Mellin integral
\[
 \mathcal MF(s)=\int_0^\infty F(x)x^{s-1}\,dx
\]
entire for $F\in\mathscr B$ and give
$\mathcal M(DF)(s)=s\mathcal MF(s)$, including the vanished endpoint
terms.

For clarity the seed and all its constants can be constructed directly.
Put
\[
 \phi_0(x)=(4\pi^2x^4-6\pi x^2)e^{-\pi x^2},\qquad
 \vartheta(z)=\sum_{n\in\mathbb Z}e^{-\pi n^2z^2},\qquad
 f_0=D(D-1)\vartheta=\Theta\phi_0.                              \tag{GC.3}
\]
Gaussian integration gives $\int4\pi^2x^4e^{-\pi x^2}\,dx=3$ and
$\int6\pi x^2e^{-\pi x^2}\,dx=3$ on $\mathbb R$, and
$\phi_0(0)=0$, so $\phi_0\in V$. The Gaussian Fourier transform
in the convention $\widehat f(y)=\int f(x)e^{-2\pi ixy}\,dx$ is
$e^{-\pi y^2}$: differentiation under the integral and integration by
parts give $\widehat f'=-2\pi y\widehat f$, and the integral at zero
is one. Periodizing this Gaussian gives an absolutely convergent
Fourier series with those computed coefficients. Evaluating it at zero
proves $\vartheta(1/x)=x\vartheta(x)$ for $x>0$. Holomorphic
continuation gives the same identity for $|\arg z|<\pi/4$, where
both theta series converge normally on compact subsets.

For $RF(x)=x^{-1}F(1/x)$, direct differentiation gives
$DR=R(1-D)$. Hence $D(D-1)$ commutes with $R$ and
\[
 f_0(1/z)=z f_0(z).                                             \tag{GC.4}
\]
On every closed narrower sector, each derivative of $f_0$ at infinity
is bounded by a polynomial times $e^{-c|z|^2}$ for some $c>0$.
Equation (GC.4), differentiated finitely many times, proves rapid
decay at zero too. In particular $f_0\in\mathscr B$.

For $\Re s>1$, absolute convergence permits termwise Mellin integration
of $2\sum_{n\ge1}\phi_0(nx)$. The change $y=\pi n^2x^2$ gives
\[
 \begin{split}
 \mathcal Mf_0(s)
 &=\pi^{-s/2}\zeta(s)
       \{4\Gamma(s/2+2)-6\Gamma(s/2+1)\}\\
 &=s(s-1)\pi^{-s/2}\Gamma(s/2)\zeta(s)=:g(s)=2\xi(s).
 \end{split}                                                   \tag{GC.5}
\]
The first line retains the theta factor two and the integration
Jacobian $1/2$. The gamma recurrence proves the second line. The
Mellin integral gives its entire continuation.

Let $h(s)=\prod_\rho(s-\rho)^{m_\rho}$ be a finite monic packet of
actual nontrivial zeros of $g$, each selected with its complete order.
The empty packet is $h=1$. Write $d=\deg h$ and $v_h=g/h$.
At each selected centre, cancelling its exact local factor constructs
the analytic value of $v_h$; hence $v_h$ is entire. It is nonzero at
each selected centre and is not identically zero.

We record the bounds that justify the integrals below. For any closed
real interval of $\sigma$, any $0<\beta<\pi/4$, and any $j\ge0$,
\[
 |v_h(\sigma+it)|\le C_{j,\beta,h}(1+|t|)^{-j}e^{-\beta|t|}.
                                                               \tag{GC.6}
\]
To prove this first for $g$, rotate its Mellin contour to
$z=e^{i\beta}x$ for $t\ge0$ and to $z=e^{-i\beta}x$ for $t<0$.
The sector estimates after (GC.4) make the two closing arcs vanish,
uniformly for the stated $\sigma$ interval. The rotation contributes
$e^{\mp\beta t}$ and a unit-modulus factor depending on $\sigma$.
Repeated integration by parts along that ray gives every displayed
power of $(1+|t|)^{-1}$; all ray endpoint terms vanish by the same
estimates. For $|t|$ larger than twice the moduli of all packet
centres relative to that $\sigma$ interval, monicity gives a lower
bound for $|h(\sigma+it)|$ by a positive constant times $|t|^d$.
On the remaining compact rectangle the cancelled entire quotient is
bounded. This proves (GC.6) for $v_h$.

For any real $\sigma$ define
\[
 F_h(x)=\frac1{2\pi}\int_{\mathbb R}
              v_h(\sigma+it)x^{-\sigma-it}\,dt.                \tag{GC.7}
\]
Horizontal rectangle estimates from (GC.6) prove independence of
$\sigma$. Differentiating in $x$ multiplies the integrand by a
power of $\sigma+it$ and remains absolutely convergent. Choosing
$\sigma$ arbitrarily large positive at infinity and arbitrarily
large negative at zero proves $F_h\in\mathscr B$. Fourier inversion
in $\log x$, with the displayed $1/(2\pi)$ constant, gives
$\mathcal MF_h=v_h$. Consequently Mellin injectivity gives
\[
 h(D)F_h=f_0=\Theta\phi_0.                                    \tag{GC.8}
\]
Define the actual density, mass, and Laplace integral by
\[
 w_h(t)=\frac{|v_h(1/2+it)|^2}{2\pi},\quad
 \mu_h=\int_{\mathbb R}w_h(t)\,dt,\quad
 M_h(\theta)=\int_{\mathbb R}e^{\theta t}w_h(t)\,dt.
                                                               \tag{GC.9}
\]
Equation (GC.6) proves holomorphy of $M_h$ on
$|\Re\theta|<\pi/2$, with
$M_h^{(j)}(\theta)=\int t^je^{\theta t}w_h(t)\,dt$ there.
Indeed any compact substrip is dominated using a $\beta$ for which
$2\beta$ is larger than its real width. The mass is positive because
$v_h$ is a nonzero analytic function; its real-line zeros are discrete.
Every original spectral moment and its $2\pi$ factor are retained.

\subsection{The gamma reference with its actual mass}

Fix $\lambda>0$ and write $\alpha=2\lambda$. Define
\[
 d\sigma_\lambda(t)=r_\lambda(t)\,dt,
 \quad r_\lambda(t)=\frac{|\Gamma(\lambda+it/2)|^2}{2\pi},
 \quad c_\lambda=2^{1-2\lambda}\Gamma(2\lambda).
                                                               \tag{GC.10}
\]
The product of two gamma integrals, changed by
$(a,b)=(ry,(1-r)y)$ with Jacobian $y$, proves the beta integral.
Its substitution $r/(1-r)=e^x$ gives
\[
 \int_{\mathbb R}(2\cosh(x/2))^{-2\lambda}e^{ivx}\,dx
 =\frac{\Gamma(\lambda+iv)\Gamma(\lambda-iv)}{\Gamma(2\lambda)}.
\]
The integrand before Fourier transformation and its second derivative
are integrable. Two integrations by parts make its transform
integrable, so Fourier inversion is valid. Alternatively, multiplying
that transform by $e^{-\varepsilon v^2}$ gives an absolutely
integrable Gaussian approximation and proves the same inversion by
dominated convergence. The substitution $t=2v$, with Jacobian two,
therefore proves
\[
 \int_{\mathbb R}e^{iyt}r_\lambda(t)\,dt
       =c_\lambda(\cosh y)^{-2\lambda}.                        \tag{GC.11}
\]
The beta integrand extends holomorphically to $|\Im x|<\pi$ using
the branch positive on $\mathbb R$. On every closed narrower strip
it is bounded by a constant times $e^{-\lambda|\Re x|}$.
Shifting the contour up for $v>0$ and down for $v<0$ proves
$|\Gamma(\lambda+iv)|^2\le C_{\lambda,b}e^{-b|v|}$ for every
$b<\pi$. Thus $\sigma_\lambda$ has exponential moments of every
order less than $\pi/2$ in the original $t$ variable.

Fourier uniqueness for finite measures and multiplication of (GC.11)
give the literal sum density
\[
 r_{\lambda,k}:=r_\lambda^{*k}
   =\frac{c_\lambda^k}{c_{k\lambda}}r_{k\lambda},\qquad
 \frac{c_\lambda^k}{c_{k\lambda}}
   =2^{k-1}\frac{\Gamma(2\lambda)^k}{\Gamma(2k\lambda)}.
                                                               \tag{GC.12}
\]
Its mass is $c_\lambda^k$. Analytic continuation justified by the
proved exponential moments gives
\[
 Z_{\lambda,k}^\Gamma(\theta)
 =\int e^{\theta u}r_{\lambda,k}(u)\,du
 =c_\lambda^k(\cos\theta)^{-k\alpha},
 \qquad |\Re\theta|<\pi/2.                                   \tag{GC.13}
\]
The power uses the analytic logarithm of $\cos\theta$ equal to zero
at zero; this cosine has no zeros on the stated strip.

For $|z|<1$, let
\[
 \begin{split}
 \mathcal G_\lambda(z,t)
  &=(1-iz)^{-\lambda+it/2}(1+iz)^{-\lambda-it/2}\\
  &=\sum_{n\ge0}\frac{b_n^{(\lambda)}(t)}{n!}z^n.
 \end{split}                                                   \tag{GC.14}
\]
Both logarithms are the principal logarithms and vanish at zero,
since $1\pm iz$ have positive real part. For real $z,w$ near zero,
the Laplace transform and the elementary cosine addition identity give
\[
 \int\mathcal G_\lambda(z,t)\mathcal G_\lambda(w,t)
              \,d\sigma_\lambda(t)
 =c_\lambda(1-zw)^{-\alpha}.
\]
On a sufficiently small closed bidisc all derivatives of the
integrands are bounded by a polynomial times $e^{a|t|}$ with
$a<\pi/2$. The exponential moment proof permits their integration.
Coefficient comparison then proves, with the monic leading
coefficient obtained from $\arctan z=z+O(z^3)$,
\[
 \begin{split}
 \int b_n^{(\lambda)}b_j^{(\lambda)}\,d\sigma_\lambda
   &=c_\lambda n!(\alpha)_n\,\delta_{nj},\\
 b_0^{(\lambda)}=1,\quad b_1^{(\lambda)}=t,\quad
 b_{n+1}^{(\lambda)}
   &=t b_n^{(\lambda)}-n(n+\alpha-1)b_{n-1}^{(\lambda)}.
 \end{split}                                                   \tag{GC.15}
\]
The recurrence follows directly by differentiating (GC.14):
$(1+z^2)\partial_z\mathcal G_\lambda=(t-\alpha z)\mathcal G_\lambda$.
The polynomials have real coefficients.

They are complete in $L^2(r_{\lambda,k}\,du)$ at parameter $k\lambda$.
In fact, an orthogonal vector $f$ defines the finite complex measure
$f(u)r_{\lambda,k}(u)\,du$. Cauchy--Schwarz and exponential
integrability give its Fourier transform on a nonempty complex strip.
All derivatives at zero vanish by polynomial orthogonality, so
analytic uniqueness and Fourier uniqueness imply $f=0$ almost
everywhere. This argument uses the full measure mass.

\subsection{The sum projection and exact coefficient powers}

On the actual spaces
\[
 \mathcal H_k=L^2(\mathbb R^k,\prod_i d\sigma_\lambda(t_i)),
 \qquad\mathcal H_\Sigma=L^2(\mathbb R,r_{\lambda,k}(u)\,du),
\]
let $U_\Sigma f(\mathbf t)=f(\sum_i t_i)$. Fubini and (GC.12)
prove that it is an isometry. The affine fibre chart has Jacobian one,
so its adjoint for $k\ge2$ is
\[
 (C_\Sigma F)(u)=\frac1{r_{\lambda,k}(u)}
 \int_{\mathbb R^{k-1}}F(t_1,\ldots,t_{k-1},u-\sum_{i<k}t_i)
 \prod_{i<k}r_\lambda(t_i)r_\lambda(u-\sum_{i<k}t_i)
 \,d^{k-1}t.                                                  \tag{GC.16}
\]
The denominator is positive everywhere. Fibre Cauchy--Schwarz and
Fubini give $\|C_\Sigma F\|\le\|F\|$, so the formula defines
the indicated adjoint on all of $\mathcal H_k$. When $k=1$, both
maps are the identity. In every case $C_\Sigma U_\Sigma=1$ and
$P_\Sigma=U_\Sigma C_\Sigma$ is the orthogonal projection. The maps
\[
 F\mapsto(C_\Sigma F,(1-P_\Sigma)F),\qquad
 (f,n)\mapsto U_\Sigma f+n,\quad n\in\ker C_\Sigma              \tag{GC.17}
\]
are inverse Hilbert isometries. This retains the entire relative fibre.

Multiplying the $k$ generating functions gives a finite identity at
each coefficient:
\[
 b_n^{(k\lambda)}(\sum_i t_i)=
 \sum_{n_1+\cdots+n_k=n}\frac{n!}{\prod_i n_i!}
                    \prod_i b_{n_i}^{(\lambda)}(t_i).
                                                               \tag{GC.18}
\]
Pairing it with a product polynomial and using (GC.15) leaves only
its total degree $n=\sum_i n_i$. The pairing at that degree is
$c_\lambda^k n!\prod_i(\alpha)_{n_i}$ and the sum norm is
$c_\lambda^k n!(k\alpha)_n$. Completeness therefore proves
\[
 C_\Sigma\prod_i b_{n_i}^{(\lambda)}
   =\frac{\prod_i(\alpha)_{n_i}}{(k\alpha)_n}b_n^{(k\lambda)}.
                                                               \tag{GC.19}
\]
The squared norm of the retained complementary polynomial is exactly
\[
 c_\lambda^k\left[
  \prod_i n_i!(\alpha)_{n_i}
  -\frac{n!\prod_i(\alpha)_{n_i}^2}{(k\alpha)_n}\right],
                                                               \tag{GC.20}
\]
by the orthogonal decomposition (GC.17).

For the original arithmetic density define, for every $\lambda>0$,
\[
 \begin{split}
 A_{h,\lambda}(t)&=\frac{v_h(1/2+it)}{\Gamma(\lambda+it/2)},
 \quad B_{h,\lambda}=|A_{h,\lambda}|^2,
 \quad w_h=B_{h,\lambda}r_\lambda,\\
 c_{h,n}&=\frac{\int b_n^{(\lambda)}(t)w_h(t)\,dt}
                 {c_\lambda n!(\alpha)_n},
 \quad\mathcal A_h(z)=\sum_{n\ge0}(\alpha)_nc_{h,n}z^n,\\
 d_{h,k,n}&=[z^n]\mathcal A_h(z)^k.
 \end{split}                                                   \tag{GC.21}
\]
The notation $\mathcal A_h$ retains the fixed parameter $\lambda$.
The complex amplitude $A_{h,\lambda}$ and the generating function
$\mathcal A_h$ are different specified maps. The integrals in
(GC.21) converge by (GC.6), before any convergence claim about the
power series. Each $d_{h,k,n}$ is the finite expression
\[
 d_{h,k,n}=\sum_{n_1+\cdots+n_k=n}
                   \prod_i(\alpha)_{n_i}c_{h,n_i}.             \tag{GC.22}
\]

Let $m_{h,k}=w_h^{*k}$. Absolute moments and Fubini give, by (GC.18),
\[
 \int_{\mathbb R}b_n^{(k\lambda)}(u)m_{h,k}(u)\,du
       =c_\lambda^k n!d_{h,k,n}.                              \tag{GC.23}
\]
This finite identity holds for every $\lambda>0$ and every $n$;
it does not require an infinite product to be interchanged with an
integral. In particular $c_{h,0}=\mu_h/c_\lambda$ and the total
mass is $\mu_h^k$.

\subsection{The supplied analytic generating join, with its branches}

\begin{theorem}[The exact analytic generating function]
For the original arithmetic measure and every $\lambda>0$, the series
in (GC.21) converges holomorphically on $|z|<1$ and is exactly
\[
 \boxed{\mathcal A_h(z)=\frac1{c_\lambda}
            (1+z^2)^{-\lambda}M_h(\arctan z).}                 \tag{GC.24}
\]
Here
$\arctan z=(\log(1+iz)-\log(1-iz))/(2i)$, and every logarithm
is the branch inherited from zero. For real $|\theta|<\pi/4$,
\[
 \boxed{M_h(\theta)=c_\lambda\sec^{2\lambda}\theta\,
                        \mathcal A_h(\tan\theta).}             \tag{GC.25}
\]
With the original tensor mass and reference mass,
\[
 \boxed{\frac{Z_{h,k}(\theta)}{Z_{\lambda,k}^\Gamma(\theta)}
       =\mathcal A_h(\tan\theta)^k,\quad
 Z_{h,k}(\theta)=\int e^{\theta u}m_{h,k}(u)\,du=M_h(\theta)^k.}
                                                               \tag{GC.26}
\]
\end{theorem}
\begin{proof}
For $|z|<1$, both $1\pm iz$ have positive real part. Their principal
logs are analytic, and their sum is the analytic logarithm of
$1+z^2$ vanishing at zero. Direct substitution in (GC.14) gives
\[
 \mathcal G_\lambda(z,t)=(1+z^2)^{-\lambda}
                              e^{t\arctan z}.                 \tag{GC.27}
\]
There is no omitted phase in this identity: the $it/2$ coefficient
multiplies $\log(1-iz)-\log(1+iz)=-2i\arctan z$.

The ratio $w=(1+iz)/(1-iz)$ has
\[
 \Re w=\frac{1-|z|^2}{|1-iz|^2}>0,\qquad
 \tan(2\Re\arctan z)=\frac{2\Re z}{1-|z|^2}.
\]
The argument of $w$ lies in $(-\pi/2,\pi/2)$ and equals
$2\Re\arctan z$. Thus for every $0\le r<1$,
\[
 |z|\le r\quad\Longrightarrow\quad
 |\Re\arctan z|\le\arctan r<\pi/4,
 \qquad
 |\mathcal G_\lambda(z,t)|
 \le(1-r^2)^{-\lambda}e^{\arctan(r)|t|}.                       \tag{GC.28}
\]
For the angle inequality use monotonicity of arctangent and
$2|\Re z|/(1-|z|^2)\le2r/(1-r^2)$, together with
$2\arctan r\in[0,\pi/2)$. For the modulus inequality use
$|1+z^2|\ge1-r^2$ and real $\lambda>0$.

Equation (GC.6) makes the right side of (GC.28) integrable against
$w_h(t)\,dt$, uniformly on every compact subdisc. Dominated
integration therefore defines a holomorphic function
$c_\lambda^{-1}\int\mathcal G_\lambda(z,t)w_h(t)\,dt$ on the
whole unit disc. One explicit justification is Morera's theorem:
domination permits exchanging the integral with each triangle
contour, and the inner holomorphic contour integral is zero.
For its derivatives on $|z|\le r$, choose $r<R<1$ and apply
Cauchy's derivative formula on circles of radius $R-r$; the
kernel derivative $|\partial_z^n\mathcal G_\lambda(z,t)|$ is bounded by
$n!(R-r)^{-n}(1-R^2)^{-\lambda}e^{\arctan(R)|t|}$.
This integrable bound justifies every derivative interchange.

At zero, (GC.14) now gives
\[
 \mathcal A_h^{(n)}(0)=\frac1{c_\lambda}
               \int b_n^{(\lambda)}(t)w_h(t)\,dt
           =n!(\alpha)_nc_{h,n}.                              \tag{GC.29}
\]
Thus the Taylor series of this holomorphic integral is exactly the
original coefficient series. Substitution of (GC.27) in the integral
proves (GC.24). All its constants occur before the integration.

For real $|\theta|<\pi/4$, $|\tan\theta|<1$,
$\arctan(\tan\theta)=\theta$ on the specified branch, and
$1+\tan^2\theta=\sec^2\theta>0$. Solving (GC.24) for
$M_h(\theta)$ gives (GC.25). Finally the absolute integral of
$\prod_i e^{\theta t_i}w_h(t_i)$ is finite throughout
$|\Re\theta|<\pi/2$. Fubini gives $Z_{h,k}=M_h^k$ with no
mass division. Divide the kth power of (GC.25) by (GC.13) to
obtain (GC.26). At zero this ratio is explicitly
$(\mu_h/c_\lambda)^k$.
\end{proof}

The inverse also has a precise complex domain. If
$|\Re\theta|<\pi/4$, then
\[
 |\tan\theta|^2=
 \frac{\cosh(2\Im\theta)-\cos(2\Re\theta)}
      {\cosh(2\Im\theta)+\cos(2\Re\theta)}<1.
\]
The denominator is positive and the branch satisfies
$\arctan(\tan\theta)=\theta$, first on a neighbourhood of zero
and then on this connected strip by analytic uniqueness. Hence
(GC.25)--(GC.26) extend to that strip with the analytic branch of
the cosine power in (GC.13). No logarithm of $\mathcal A_h$ is
asserted: it can have complex zeros, whereas all real-tilt integrals
$M_h(\theta)$ are strictly positive.

For each finite derivative there is also an exact moment formula.
Write $b_n^{(\lambda)}(t)=\sum_{j=0}^n b_{n,j}t^j$. Then
\[
 n!(\alpha)_nc_{h,n}
  =\frac1{c_\lambda}\sum_{j=0}^n b_{n,j}M_h^{(j)}(0),\qquad
 d_{h,k,n}=\frac1{n!}
      \left.\frac{d^n}{dz^n}\mathcal A_h(z)^k\right|_{z=0}.
                                                               \tag{GC.30}
\]
The first identity follows by the finite polynomial expansion and
(GC.9); the second is Taylor's formula, whose Leibniz sum is
exactly (GC.22). The triangular coefficients $b_{n,n}=1$ give an
invertible finite map between moments through degree n and coefficients
through degree n, with the full diagonal factors
$c_\lambda n!(\alpha)_n$ retained.

\subsection{The complete density expansion for the original bounded references}

The finite identities and analytic join above hold for every
$\lambda>0$. For the original full-quartet reference choose
$\lambda=1/4$. More generally choose the actual shifted reference
\[
 \lambda=\tfrac14+L,\qquad
 L\text{ an integer},\quad L\ge\max(0,3-d).                    \tag{GC.31}
\]
These choices admit the following proved bounded multiplier and hence
the complete $L^2$ density expansion.

Indeed at $s=1/2+it$ the defining g identity gives
\[
 A_{h,1/4}(t)=
 \frac{-\pi^{-1/4}(t^2+1/4)e^{-it\log\pi/2}
                  \zeta(1/2+it)}{h(1/2+it)}.                  \tag{GC.32}
\]
This expression is continued through each selected zero by the
already defined entire quotient. Integrating the step function
$\lfloor x\rfloor$, initially for $\Re s>1$, gives
$\zeta(s)=s/(s-1)-s\int_1^\infty\{x\}x^{-s-1}\,dx$.
The last integral is holomorphic for $\Re s>0$, so the identity
continues there away from 1. On the original line it gives
$|\zeta(s)|\le1+2\sqrt{t^2+1/4}\le2(1+|t|)$.
For a nonempty packet put
$T_h=\max(1,2\max_\rho|\rho-1/2|)$; for $h=1$ put $T_h=1$.
For $|t|\ge T_h$, every original packet factor has modulus at
least $|t|/2$. Gamma recurrence then gives
\[
 \begin{split}
 A_{h,1/4+L}(t)&=
   \frac{A_{h,1/4}(t)}{\prod_{j=0}^{L-1}(j+1/4+it/2)},\\
 B_{h,1/4+L}(t)&\le
       \frac{25\,2^{2(d+L)}}{\sqrt\pi}|t|^{6-2d-2L}.
 \end{split}                                                   \tag{GC.33}
\]
For the constant 25, the two squared numerator estimates at $|t|\ge1$
are $(t^2+1/4)^2\le25t^4/16$ and
$|\zeta(s)|^2\le16t^2$. Each shift factor has modulus at least
$|t|/2$. The exponent in (GC.33) is nonpositive for (GC.31).
Thus a completely specified finite bound is
\[
 C_{h,L}=\max\left\{
 \sup_{|t|\le T_h}B_{h,1/4+L}(t),
 \frac{25\,2^{2(d+L)}}{\sqrt\pi}T_h^{6-2d-2L}\right\}>0.
                                                               \tag{GC.34}
\]
The compact term is finite by continuity of the cancelled multiplier.
This is a mathematical bound with a specified supremum, not a claim
that a numerical packet supremum has already been enclosed.

For these references $B\in L^2(d\sigma_\lambda)$, because the
reference has finite mass and $0\le B\le C_{h,L}$. Fibre integration
of $w_h=Br_\lambda$ proves
\[
 B_{h,\lambda;k}=C_\Sigma(B^{\otimes k}),\qquad
 m_{h,k}=r_{\lambda,k}B_{h,\lambda;k},\qquad
 0\le B_{h,\lambda;k}\le C_{h,L}^k.                            \tag{GC.35}
\]
For $k\ge2$ it is positive at every sum coordinate: the zeros of B
form a countable discrete set, and their forbidden fibre coordinates
are a countable union of affine hyperplanes of measure zero. Every
remaining reference factor is positive. For $k=1$ the original
isolated zeros remain.

Combining (GC.23), completeness, and the sum norm in (GC.15) gives
the full, convergent $L^2(r_{\lambda,k}\,du)$ identity
\[
 B_{h,\lambda;k}(u)=\sum_{n\ge0}
                \frac{d_{h,k,n}}{(k\alpha)_n}b_n^{(k\lambda)}(u).
                                                               \tag{GC.36}
\]
In particular the norm of the retained relative density vector is
\[
 \left\|B^{\otimes k}-U_\Sigma B_{h,\lambda;k}\right\|^2
   =\left(\int B^2d\sigma_\lambda\right)^k
      -c_\lambda^k\sum_{n\ge0}
                       \frac{n!}{(k\alpha)_n}|d_{h,k,n}|^2.
                                                               \tag{GC.37}
\]
This is Pythagoras for (GC.17), so every norm and mass factor is
retained. A truncation of (GC.36) is not asserted pointwise positive.

\subsection{Original-coordinate finite Grams and every tilt derivative}

Let $\mathcal P_N=\mathbb C[S]_{\le N}$ for $N\ge0$, with original
ordered basis $1,S,\ldots,S^N$, and set $c=k/2$. The arithmetic
source Gram and its reference are
\[
 \begin{split}
 (\mathsf M_{h,k,N}(\theta))_{ij}
   &=\int(c-iu)^i(c+iu)^j e^{\theta u}m_{h,k}(u)\,du,\\
 (\mathsf M_{\lambda,k,N}^\Gamma(\theta))_{ij}
   &=\int(c-iu)^i(c+iu)^j e^{\theta u}r_{\lambda,k}(u)\,du,
 \qquad 0\le i,j\le N. 
 \end{split}                                                   \tag{GC.38}
\]
They are holomorphic entrywise for $|\Re\theta|<\pi/2$ and
positive definite at every real tilt in that strip: a nonzero
polynomial has only finitely many real-line roots and each measure
is positive almost everywhere. Let $\mathsf H_{h,k,N}$ denote
the corresponding real-u monomial Hankel matrix. The exact matrix
of the unchanged coordinate substitution is
\[
 (\mathsf L_{N+1})_{aj}=
 \begin{cases}\binom ja c^{j-a}i^a,&0\le a\le j\le N,\\
 0,&a>j,
 \end{cases}\qquad
 \det\mathsf L_{N+1}=i^{N(N+1)/2}.                            \tag{GC.39}
\]
Expanding each original S monomial proves
\[
 \mathsf M_{h,k,N}=\mathsf L_{N+1}^*\mathsf H_{h,k,N}
                                     \mathsf L_{N+1}            \tag{GC.40}
\]
and the identical congruence for the reference. Since the displayed
determinant has modulus one, it also proves exact equality of the
two coordinate determinants. The congruence, rather than a relabeling
of a matrix, preserves every complex phase in its entries.

For an integer $r\ge0$, expand the specified polynomial uniquely as
\[
 u^r(c-iu)^i(c+iu)^j
 =\sum_{n=0}^{i+j+r}e_{ijr,n}^{(k,\lambda)}b_n^{(k\lambda)}(u).
                                                               \tag{GC.41}
\]
The monic triangular basis makes these coefficients finite, unique,
and computable by polynomial subtraction in the written order.
Differentiation of (GC.38) and (GC.23) now give the exact formula
\[
 \boxed{\left.\partial_\theta^r
      (\mathsf M_{h,k,N})_{ij}\right|_{\theta=0}
 =c_\lambda^k\sum_{n=0}^{i+j+r}
                 e_{ijr,n}^{(k,\lambda)}n!d_{h,k,n}.}           \tag{GC.42}
\]
The derivative interchange is justified by the exponential bounds
after (GC.9), or directly by Fubini on the k-fold product. Since
$d_{h,k,n}$ uses exactly the finite coefficient power (GC.22), the
entire rth derivative matrix is determined by
$c_{h,0},\ldots,c_{h,2N+r}$. At zero tilt $r=0$, the list ends at
$2N$. There is no truncation error in this finite statement.

Equation (GC.42) supplies each derivative at zero from its stated
finite coefficient list. Evaluation at a nonzero tilt is related to
those derivatives by the following exact analytic Taylor identity,
valid for $|\theta|<\pi/2$:
\[
 (\mathsf M_{h,k,N}(\theta))_{ij}
 =\sum_{r\ge0}\frac{\theta^r}{r!}
       c_\lambda^k\sum_{n=0}^{i+j+r}
                    e_{ijr,n}^{(k,\lambda)}n!d_{h,k,n}.        \tag{GC.43}
\]
Absolute convergence follows by bounding the absolute exponential
series by $e^{|\theta||u|}$ times the original polynomial modulus.
For larger imaginary components within the full holomorphic strip,
analytic continuation of these entries is the defined original
Laplace integral (GC.38). This specifies the relation between all
finite coefficient jets and the full tilted metric.

\subsection{The original cyclic quotient and the exact cochain correction}

For a nonempty packet let $E_h=\mathbb C[s]/(h)$ and
$\mathcal E_h=E_h^{\otimes k}$. Let $a$ be the element
$s_1+\cdots+s_k$ of this finite algebra. The kernel of
$\mathbb C[S]\to\mathcal E_h$, $P\mapsto P(a)$, is generated by
its unique monic polynomial $\chi$. It is nonzero since the powers
of $a$ lie in a finite-dimensional vector space. Put
$q=\deg\chi\ge1$, $C=\mathbb C[S]/(\chi)$, and $A=M_S:C\to C$.
The induced map $C\to\mathcal E_h$ is injective by the definition
of the kernel. For each selected zero, the complete Taylor value of
$g/h$ is nonzero at its centre; its residue
$\upsilon_h=j_h(g/h)\in E_h$ is consequently a unit in every
local factor and in $E_h$. Thus the original arithmetic injection
\[
 \eta_{h,k}:C\hookrightarrow\mathcal E_h,\qquad
 [P]\mapsto\upsilon_h^{\otimes k}P(a)                           \tag{GC.44}
\]
retains the entire local unit and every selected multiplicity.

The original analytic polynomial map is
\[
 \mathcal V_{h,k}P=P(D_1+\cdots+D_k)F_h^{\otimes k}.             \tag{GC.45}
\]
Its tensor Mellin transform is
$P(s_1+\cdots+s_k)\prod_i v_h(s_i)$. Taking complete selected
jets gives precisely (GC.44) after the quotient $\pi_\chi$.
No modulus operation has been applied to that arithmetic jet map.

Successive monic divisions, in the order $s_1,\ldots,s_k$, provide
specific polynomials $Q_i(s_1,\ldots,s_k)$ with
\[
 \chi(s_1+\cdots+s_k)=\sum_{i=1}^k h(s_i)Q_i(s_1,\ldots,s_k).
                                                               \tag{GC.46}
\]
To prove this, divide in $s_1$, retain its remainder, divide that
remainder in $s_2$, and continue. The final remainder has degree
less than d in every coordinate. Its residue in $\mathcal E_h$ is
zero by the defining property of $\chi$, and the tensor remainder
monomials form a basis, so every coefficient is zero. This proves
the displayed identity and fixes the division order.

The operators $D_i$ commute on the tensor complex. For a polynomial
$P$, define the actual cochain of degree $k-1$ by
\[
 \begin{split}
 \mathcal P_{\chi,h,k}(P)=
 \sum_{i=1}^k(-1)^{i-1}Q_i(D_1,\ldots,D_k)
 P(D_1+\cdots+D_k)\quad\quad\quad\quad\quad\quad\quad\\[-4pt]
 \cdot\bigl(F_h^{\otimes(i-1)}\otimes\phi_0
                                \otimes F_h^{\otimes(k-i)}\bigr).
 \end{split}                                                   \tag{GC.47}
\]
Every summand belongs to the indicated tensor source, since D
preserves V and $\mathscr B$. The first $i-1$ factors have cochain
degree one. The tensor differential on the ith factor contributes
$(-1)^{i-1}$, cancelling the written sign. The other factors have
zero differential. Equations (GC.8) and (GC.46) therefore give
\[
 d_{\rm tot}\mathcal P_{\chi,h,k}(P)
       =\mathcal V_{h,k}(\chi P).                             \tag{GC.48}
\]
This identifies the cochain primitive, its sign, and its boundary in
their different cochain degrees.

For the full quotient at finite source degree take $N\ge q-1$ and
$m=N-q+1$. With $\mathcal P_{-1}=0$, monic division gives the
exact sequence in original coordinates
\[
 0\longrightarrow\mathcal P_{m-1}
 \xrightarrow{M_\chi}\mathcal P_N
 \xrightarrow{J_N}C\longrightarrow0,
 \quad M_\chi P=\chi P,\quad J_NP=[P]_\chi.                    \tag{GC.49}
\]
The kernel assertion is divisibility, and the first q monomial
columns prove surjectivity. Let $B_\chi$ be the $(N+1)\times m$
matrix of $M_\chi$ in the original ordered bases. Its entries are
$[S^i](\chi(S)S^j)$; it is independent of the metric and tilt.
The relation Gram and every zero-tilt derivative are therefore
\[
 \mathsf B_{h,k,N}(\theta)=B_\chi^*\mathsf M_{h,k,N}(\theta)B_\chi,
 \qquad
 \mathsf B_{h,k,N}^{(r)}(0)=B_\chi^*\mathsf M_{h,k,N}^{(r)}(0)B_\chi.
                                                               \tag{GC.50}
\]
Equivalently its polynomial integrand is
$u^r|\chi(c+iu)|^2(c-iu)^i(c+iu)^j$, where $0\le i,j<m$.
Its degree is at most $r+2q+2(m-1)=r+2N$, so the relation derivative
uses the same list through $2N+r$ as the source derivative. Both
complex polynomial factors, their conjugates, and all $\chi$ coefficients
remain in this formula.

For a fixed real admitted tilt write $M_h=\mathsf M_{h,k,N}(\theta)$
and $M_\Gamma=\mathsf M_{\lambda,k,N}^\Gamma(\theta)$.
Their canonical quotient Grams and minimum lifts are
\[
 G_h=(J_NM_h^{-1}J_N^*)^{-1},\quad
 R_h=M_h^{-1}J_N^*G_h,
\qquad
 G_\Gamma=(J_NM_\Gamma^{-1}J_N^*)^{-1},\quad
 R_\Gamma=M_\Gamma^{-1}J_N^*G_\Gamma.                           \tag{GC.51}
\]
Surjectivity of J and positivity of the source Grams prove every
inverse exists. These formulas give $JR_h=JR_\Gamma=I_C$ and
make each lift orthogonal to $\ker J$ in its own metric, proving
its minimum property by the orthogonal Pythagorean identity.
For $m>0$ put
\[
 K=(B_\chi^*M_hB_\chi)^{-1}B_\chi^*M_hR_\Gamma:
                         C\longrightarrow\mathbb C^m.
                                                               \tag{GC.52}
\]
The relation Gram is positive definite because $B_\chi$ is
injective. The section $R_\Gamma-B_\chi K$ still has quotient
identity and is orthogonal to the relation space in $M_h$.
The preceding uniqueness gives the exact section map
\[
 \boxed{R_h=R_\Gamma-B_\chi K.}                               \tag{GC.53}
\]
At $N=q-1$, $m=0$, take K to be the zero map to the zero vector
space and omit the inverse in (GC.52); (GC.53) remains the identity
of the unique degree-less-than-q representatives.

For $u_0\in C$, let $\kappa_{u_0}(S)$ be the polynomial whose
coefficient vector is $Ku_0$. Combining (GC.48) and (GC.53) proves
the exact cochain correction
\[
 \boxed{\mathcal V_{h,k}(R_hu_0-R_\Gamma u_0)
  =-\mathcal V_{h,k}(\chi\kappa_{u_0})
  =d_{\rm tot}\bigl(-\mathcal P_{\chi,h,k}(\kappa_{u_0})\bigr).}
                                                               \tag{GC.54}
\]
Thus applying $\mathcal V$ gives the boundary in degree k; the
explicit negative cochain on the right is its primitive in degree
$k-1$. Passing to the original arithmetic quotient sends this
boundary to zero. Explicitly, for the original set of labels
$\mathcal L$ and vector spaces $E_\ell$, write the supplied carrier
as $\{\tau\}\sqcup\bigsqcup_{\ell\in\mathcal L}(\{\ell\}\times E_\ell)$.
A family of linear maps $L_\ell:E_\ell\to F_\ell$ has the exact
support lift $\tau\mapsto\tau$ and
$(\ell,v)\mapsto(\ell,L_\ell v)$, with
$e_\ell=(\ell,0)$ in each target fibre. Substitution proves that
these lifts preserve identities and composition. Applied to the
original quotient, the relation boundary in (GC.54) therefore has
image $e_\ell$, while the external $\tau$ has image $\tau$.
The source vector and its primitive remain available before that
quotient.

\subsection{Determinant indices and the original Toda substitution}

For real admitted $\theta$, write $\mathfrak D_n$ for the original
degree-less-than-n source Gram determinant and $\mathfrak B_m$ for
the relation Gram determinant on m multiplication columns. Empty
determinants are one. The analogous reference determinants carry
the superscript $\Gamma$. For $N\ge q-1$, set
\[
 V_N=\det G_h,\quad V_N^\Gamma=\det G_\Gamma,
 \quad X_n=\mathfrak D_n/\mathfrak D_n^\Gamma,
 \quad Y_m=\mathfrak B_m/\mathfrak B_m^\Gamma.
\]
The quotient-first coefficient matrix with columns
$1,S,\ldots,S^{q-1},\chi,\chi S,\ldots,\chi S^{m-1}$
has determinant one: its columns have successively distinct degrees
and all their leading coefficients are one. Replacing its first q
columns by the corresponding minimum sections only subtracts
relation columns, so it still has determinant one. Orthogonality
then gives the exact block determinant formula
\[
 V_N=\frac{\mathfrak D_{N+1}}{\mathfrak B_{N-q+1}}
     =V_N^\Gamma T_N,\qquad
 T_N=\frac{X_{N+1}}{Y_{N-q+1}},\quad N\ge q-1.                  \tag{GC.55}
\]
In particular the first full quotient has denominator
$\mathfrak B_0=1$ and requires no relation inverse.

For completeness the reference source determinant is
\[
 \mathfrak D_n^\Gamma(\theta)=c_\lambda^{kn}
  \prod_{j=0}^{n-1}j!(k\alpha)_j
             (\cos\theta)^{-n(k\alpha+n-1)},
 \qquad
 a_n^\Gamma(\theta)=n(n+k\alpha-1)\sec^2\theta.                \tag{GC.56}
\]
Here $a_n$ is the ratio of consecutive monic squared norms in the
real-u coordinate; the norm-preserving phase of (GC.39) gives the
same ratios in the original S coordinate. A direct proof of the
Hankel identity used to obtain (GC.56) is as follows. Let $p_j$ be
the monic orthogonal polynomials for an $e^{\theta u}$-tilted
positive measure, let $h_j$ be their squared norms, and write
$up_j=p_{j+1}+b_jp_j+a_jp_{j-1}$, with $p_{-1}=0$ and $a_0=0$.
For $j\ge1$, pairing with $p_{j-1}$ gives
$a_j=h_j/h_{j-1}$, and pairing with $p_j$ gives $h_j'=b_jh_j$
for every $j\ge0$. Differentiating orthogonality for $j\ge1$ shows
$\partial_\theta p_j=-a_jp_{j-1}$: its pairings with degrees at
most $j-2$ vanish, and the remaining pairing is $-h_j$.
For $j\ge1$, write $p_j=u^j+c_ju^{j-1}+\cdots$, and set
$c_0=0$ for $p_0=1$. The derivative statement for $p_0$ is
$\partial_\theta p_0=0$. Thus $c_j'=-a_j$ and
$b_j=c_j-c_{j+1}$ hold for every $j\ge0$, and consequently
$b_j'=a_{j+1}-a_j$.
Since $\mathfrak D_n=\prod_{j<n}h_j$,
\[
 (\log\mathfrak D_n)''=\sum_{j<n}b_j'=a_n
   =\mathfrak D_{n+1}\mathfrak D_{n-1}/\mathfrak D_n^2,
   \qquad n\ge1.
\]
For the reference, the determinant values at n=0 and n=1 are 1
and (GC.13). The right side of the proposed formula (GC.56) has
second logarithmic derivative $n(k\alpha+n-1)\sec^2\theta$.
Its consecutive ratio has exactly the same value, with constants
$n!(k\alpha)_n/((n-1)!(k\alpha)_{n-1})$ and cosine exponent
$-2$. Positive induction in the displayed determinant identity
proves (GC.56) for every n.

The reference relation determinant is the determinant of the
moment matrix with entries
\[
 \partial_\theta^{i+j}\left\{
  \overline\chi(c-i\partial_\theta)\chi(c+i\partial_\theta)
                  Z_{\lambda,k}^\Gamma(\theta)\right\},
       \qquad0\le i,j<m.                                     \tag{GC.57}
\]
Differentiation under the integral gives the unchanged weight
$|\chi(c+iu)|^2r_{\lambda,k}(u)$ in this u-coordinate matrix.
The relation version of (GC.39)--(GC.40) has determinant of modulus
one, proving that (GC.57) is the original S-coordinate relation
determinant as stated. It contains the actual $\chi$ with all coefficients.

At $\theta=0$ put
\[
 Q_N=\frac{X_{N+2}X_N}{X_{N+1}^2},\qquad
 \mathcal R_N^\Gamma=\frac{V_{N-1}^\Gamma}{V_{N+1}^\Gamma}.
\]
The consecutive full-quotient formulas have domain $N\ge q$ and
give, by literal substitution in (GC.55)--(GC.56),
\[
 a_{N+1}=(N+1)(N+k\alpha)Q_N,\qquad
 \frac{V_{N-1}}{V_{N+1}}
     =\mathcal R_N^\Gamma\frac{T_{N-1}}{T_{N+1}}.
                                                               \tag{GC.58}
\]
The domain $N\ge q$ is necessary because $V_{N-1}$ occurs. The
separate original endpoint at $N=q-1$ is not replaced by a
q-dimensional metric at degree $q-2$. At every fixed admissible
real tilt the phase derivative remains the exact equality
\[
 \partial_\theta\log V_N
   =\partial_\theta\log V_N^\Gamma+
                              \partial_\theta\log T_N,
                       \qquad N\ge q-1.                       \tag{GC.59}
\]
All these derivatives are of positive real-tilt determinants. Formula
(GC.42) gives their underlying source and relation entry derivatives
with the precise finite coefficient indices; their quotient and
determinant operations use those original matrices.

For $h=1$ the finite arithmetic quotient is zero and $\chi=1$.
Its relation and source spaces coincide at every degree N, so
$\mathfrak B_{N+1}=\mathfrak D_{N+1}$,
$V_N=V_N^\Gamma=T_N=1$. The quotient lift has zero-dimensional
domain and the cochain correction (GC.54) is the zero map on that
domain. The nonzero analytic density, its mass $\mu_1^k$, its
coefficient series and (GC.24)--(GC.43) remain as constructed.
The support lift is $G(0)=\{\tau,e\}$; it does not produce a
nonzero finite spectral quotient from the analytic seed.


\endgroup
